\begin{Example}
\label{ex:fourier-heaviside-regularization}
Fourier transforms which cannot be defined classically as an integral can often be determined as limits of regularizations. As an example, we consider the Heaviside function $H(x)$. Clearly, we cannot define the Fourier transform as
\begin{equation}\label{eq:heaviside-nonconvergent-fourier-integral}
\dfrac1{\sqrt{2\pi}}\int_0^\infty e^{-i\xi x}\,\mathrm{d}x.
\end{equation}
Observe, however, that
\begin{equation}\label{eq:heaviside-exponential-regularization}
H(x)=\lim_{\epsilon\to0^+}H(x)e^{-\epsilon x}
\end{equation}
in the sense of tempered distributions, and consequently
\begin{equation}\label{eq:heaviside-fourier-regularized-limit}
\mathcal{F}[H](\xi)=\dfrac1{\sqrt{2\pi}}\lim_{\epsilon\to0^+}\int_0^\infty e^{-\epsilon x-i\xi x}\,\mathrm{d}x=\lim_{\epsilon\to0^+}\dfrac1{\sqrt{2\pi}}\dfrac1{\epsilon+i\xi}.
\end{equation}
We can evaluate this limit a little more explicitly by applying it to a test function. For any $\delta>0$, we have
\begin{equation}\label{eq:heaviside-fourier-principal-value-limit}
\begin{aligned}
\lim_{\epsilon\to0^+}\int_{-\infty}^\infty\dfrac{\phi(\xi)}{\epsilon+i\xi}\,\mathrm{d}\xi
&=\int_{|\xi|>\delta}\dfrac{\phi(\xi)}{i\xi}\,\mathrm{d}\xi
+\int_{-\delta}^{\delta}\dfrac{\phi(\xi)-\phi(0)}{i\xi}\,\mathrm{d}\xi\\
&\quad+\lim_{\epsilon\to0^+}\int_{-\delta}^{\delta}\dfrac{\phi(0)}{\epsilon+i\xi}\,\mathrm{d}\xi\\
&=\int_{|\xi|>\delta}\dfrac{\phi(\xi)}{i\xi}\,\mathrm{d}\xi
+\int_{-\delta}^{\delta}\dfrac{\phi(\xi)-\phi(0)}{i\xi}\,\mathrm{d}\xi+\pi\phi(0).
\end{aligned}
\end{equation}
We conclude that
\begin{equation}\label{eq:heaviside-fourier-principal-value-formula}
\mathcal{F}[H]=\dfrac1{\sqrt{2\pi}}\left(\dfrac1{i\xi}+\pi\delta\right),
\end{equation}
where $1/(i\xi)$ is interpreted as a principal value.
\end{Example}

\begin{Example}
\label{ex:fourier-polynomial-growth-truncation}
Let $f$ be any continuous function which has polynomial growth at infinity. Then, in the sense of tempered distributions, $f$ is the limit as $M\to\infty$ of
\begin{equation}\label{eq:polynomial-growth-function-truncation}
f_M(\bm{x})=\begin{cases}f(\bm{x}),&|\bm{x}|\leq M\\0,&|\bm{x}|>M.\end{cases}
\end{equation}
As a consequence, we find that, in the sense of tempered distributions,
\begin{equation}\label{eq:polynomial-growth-fourier-truncation-limit}
\widehat f(\bm{\xi})=(2\pi)^{-m/2}\lim_{M\to\infty}\int_{|\bm{x}|\leq M}f(\bm{x})e^{-i\bm{\xi}\cdot\bm{x}}\,\mathrm{d}\bm{x}.
\end{equation}
In particular, if $f$ is integrable at infinity, the Fourier transform of $f$ as a distribution agrees with the ordinary Fourier transform. Another way to evaluate the Fourier transform of functions with polynomial growth is therefore to approximate them by integrable functions, such as $f(\bm{x})\exp(-\epsilon|\bm{x}|^2)$. See Problem~\ref{ex:fourier-polynomial-growth-regularizations} for examples.
\end{Example}

\begin{Example}
\label{ex:fourier-uniform-sphere-mass}
Let $\delta(r-a)$ represent a uniform mass distribution on the sphere of radius $a$, i.e.,
\begin{equation}\label{eq:uniform-sphere-mass-pairing}
(\delta(r-a),\phi)=\int_{|\bm{x}|=a}\phi(\bm{x})\,\mathrm{d}S.
\end{equation}
(Of course, this is not consistent with our previous use of "$\delta$" as a distribution on $\mathbb{R}^m$, but it is a standard abuse of notation with which the reader should become accustomed.) Then the Fourier transform of $\delta(r-a)$ is given by \eqref{eq:compact-distribution-fourier-entire-formula}
\begin{equation}\label{eq:uniform-sphere-mass-fourier-integral}
\mathcal{F}[\delta(r-a)](\bm{\xi})=(2\pi)^{-m/2}\int_{|\bm{x}|=a}e^{-i\bm{\xi}\cdot\bm{x}}\,\mathrm{d}S.
\end{equation}
We want to evaluate this expression for $m=3$. We use polar coordinates with the axis aligned with the direction of $\bm{\xi}$ so that $\bm{\xi}\cdot\bm{x}=a|\bm{\xi}|\cos\theta$; we shall use $\rho$ to denote $|\bm{\xi}|$. We thus find
\begin{equation}\label{eq:uniform-sphere-mass-three-dimensional-fourier}
\begin{aligned}
\mathcal{F}[\delta(r-a)](\bm{\xi})
&=(2\pi)^{-3/2}a^2\int_0^\pi\int_0^{2\pi}e^{-ia\rho\cos\theta}\sin\theta\,\mathrm{d}\phi\,\mathrm{d}\theta\\
&=\sqrt{\dfrac2\pi}a\dfrac{\sin a\rho}\rho.
\end{aligned}
\end{equation}

\end{Example}

\begin{Example}
\label{ex:fourier-direct-product}
The Fourier transform of a direct product is the direct product of the Fourier transforms. To show this, it suffices to prove agreement for a dense set of test functions. We have
\begin{equation}\label{eq:fourier-direct-product-separable-proof}
\begin{aligned}
(\widehat f(\bm{\xi})\widehat g(\bm{\eta}),\widehat\phi(\bm{\xi})\widehat\psi(\bm{\eta}))
&=(\widehat f,\widehat\phi)(\widehat g,\widehat\psi)\\
&=(f,\phi)(g,\psi)\\
&=(f(\bm{x})g(\bm{y}),\phi(\bm{x})\psi(\bm{y})).
\end{aligned}
\end{equation}

\end{Example}
